\clearpage
\VerdierPhysicalPageAnchor{0268}%
\begingroup
\small
\setlength{\parindent}{0pt}%
\setlength{\parskip}{0pt}%
\setlength{\emergencystretch}{3em}%
\newcommand{\VerdierContentsRow}[3]{%
  \par\noindent
  \hangindent=#1\relax\hangafter=1
  \hspace*{#1}#2\nobreak
  \leaders\hbox{\kern.32em.\kern.32em}\hfill
  \nobreak\hbox{\kern.65em #3}\par}
\newcommand{\VerdierContentsChapter}[1]{%
  \par\addvspace{.55\baselineskip}\noindent\textbf{#1}\par\smallskip}
\newcommand{\VerdierContentsSection}[2]{%
  \par\addvspace{.45\baselineskip}\noindent
  \hfill\textbf{#1.}\enspace #2\hfill\mbox{}\par\smallskip}

\vspace*{4.25\baselineskip}
\VerdierSemanticAnchor{toc-title}%
\begin{center}
  {\Large\bfseries 目次}
\end{center}
\vspace{2.5\baselineskip}

\VerdierSemanticAnchor{toc-introduction}%
\VerdierContentsRow{0pt}{\textbf{序論}}{1}
\VerdierSemanticAnchor{toc-acknowledgments}%
\VerdierContentsRow{0pt}{\textbf{謝辞}}{15}

\VerdierSemanticAnchor{toc-ch1}%
\VerdierContentsChapter{第I章：加法圏における複体の圏}
\VerdierSemanticAnchor{toc-ch1-sec1}%
\VerdierContentsSection{1}{次数付き圏}
\VerdierSemanticAnchor{toc-ch1-sec1-1}%
\VerdierContentsRow{.5em}{\textbf{1.1.}\enspace $G$-次数付き圏}{17}
\VerdierSemanticAnchor{toc-ch1-sec1-2}%
\VerdierContentsRow{.5em}{\textbf{1.2.}\enspace 同型を除いて定まる圏への群作用}{19}
\VerdierSemanticAnchor{toc-ch1-sec1-3}%
\VerdierContentsRow{.5em}{\textbf{1.3.}\enspace 二つの言語の同値性}{22}
\VerdierSemanticAnchor{toc-ch1-sec1-4}%
\VerdierContentsRow{.5em}{\textbf{1.4.}\enspace 群の変更}{27}
\VerdierSemanticAnchor{toc-ch1-sec1-5}%
\VerdierContentsRow{.5em}{\textbf{1.5.}\enspace $\mathbb Z^n$-圏}{34}
\VerdierSemanticAnchor{toc-ch1-sec1-6}%
\VerdierContentsRow{.5em}{\textbf{1.6.}\enspace Koszul 余サイクルと符号規則}{43}

\VerdierSemanticAnchor{toc-ch1-sec2}%
\VerdierContentsSection{2}{加法圏における複体}
\VerdierSemanticAnchor{toc-ch1-sec2-1}%
\VerdierContentsRow{.5em}{\textbf{2.1.}\enspace $n$-重複複体}{49}
\VerdierSemanticAnchor{toc-ch1-sec2-2}%
\VerdierContentsRow{.5em}{\textbf{2.2.}\enspace 付随する単純複体}{50}
\VerdierSemanticAnchor{toc-ch1-sec2-3}%
\VerdierContentsRow{.5em}{\textbf{2.3.}\enspace 複体の圏上の次数付き構造}{54}
\VerdierSemanticAnchor{toc-ch1-sec2-4}%
\VerdierContentsRow{.5em}{\textbf{2.4.}\enspace 関手の複体への拡張}{56}
\VerdierSemanticAnchor{toc-ch1-sec2-5}%
\VerdierContentsRow{.5em}{\textbf{2.5.}\enspace ホモトピー}{60}

\VerdierSemanticAnchor{toc-ch1-sec3}%
\VerdierContentsSection{3}{完全三角}
\VerdierSemanticAnchor{toc-ch1-sec3-1}%
\VerdierContentsRow{.5em}{\textbf{3.1.}\enspace 射の写像錐}{65}
\VerdierSemanticAnchor{toc-ch1-sec3-2}%
\VerdierContentsRow{.5em}{\textbf{3.2.}\enspace 写像錐関手の性質}{69}
\VerdierSemanticAnchor{toc-ch1-sec3-3}%
\VerdierContentsRow{.5em}{\textbf{3.3.}\enspace 完全三角の性質}{78}
\VerdierSemanticAnchor{toc-ch1-sec3-4}%
\VerdierContentsRow{.5em}{\textbf{3.4.}\enspace 完全三角の関手的性質}{84}
\VerdierSemanticAnchor{toc-ch1-appendix}%
\VerdierContentsRow{0pt}{\textbf{付録：符号の選択に関する注記}}{91}
\endgroup
